% ─────────────────────────────────────────────────────────────── 1 BÀI TOÁN
\section{Bài toán và câu hỏi}

\begin{frame}{Đọc hiểu trích xuất: chỉ ra đáp án \emph{trong} văn bản --- hoặc từ chối}
\begin{columns}[T]
\begin{column}{0.52\textwidth}
\small
\textbf{Đầu vào:} đoạn văn Wikipedia tiếng Việt + câu hỏi\\
\textbf{Đầu ra:} một chuỗi \emph{con} của đoạn văn, hoặc $\varnothing$

\vspace{0.6em}
\begin{block}{Ví dụ}
\footnotesize \emph{``\dots\ \textcolor{phobert}{\textbf{Hà Nội}} là thủ đô của nước Việt Nam \dots''}\\[2pt]
Hỏi: Thủ đô của Việt Nam là gì? $\;\to\;$ \textbf{Hà Nội}\\
Hỏi: GDP của Hà Nội năm 2030? $\;\to\;$ $\varnothing$ \textcolor{muted}{(không có đáp án)}
\end{block}

\vspace{0.4em}
UIT-ViQuAD~2.0: \PctValImp\% câu validation \textbf{không có đáp án} $\Rightarrow$ mô hình phải biết \emph{từ chối}.
\end{column}
\begin{column}{0.44\textwidth}
\small
\textbf{Hai họ mô hình khác nhau đúng ở khâu tách từ} (Buổi 3)
\vspace{0.4em}

\textcolor{phobert}{\textbf{PhoBERT}}: đơn ngữ, đầu vào \emph{đã tách từ}\\
\quad{\footnotesize\code{Hà\_Nội | là | thủ\_đô}}\\[4pt]
\textcolor{xlmr}{\textbf{XLM-R}}: đa ngữ, \emph{subword} trên văn bản thô\\
\quad{\footnotesize\code{\_Hà | \_Nội | \_là | \_thủ | \_đô}}

\vspace{0.6em}
\alert{Mô hình word-level không thể trả lời \emph{nửa từ}.}\\
\footnotesize Nếu đáp án cắt ngang ``Hà\_Nội'', PhoBERT không có cách nào khớp chính xác.
\end{column}
\end{columns}
\end{frame}

\begin{frame}{Bốn câu hỏi nghiên cứu}
\small
\renewcommand{\arraystretch}{1.35}
\begin{tabularx}{\textwidth}{@{}l X X@{}}
\toprule
& \textbf{Câu hỏi} & \textbf{Đo bằng} \\
\midrule
\textbf{RQ1} & PhoBERT (word-level) có nhạy với \textbf{ranh giới từ} hơn XLM-R (subword)? & trần EM do biên từ · loại lỗi · so sánh cặp McNemar \\
\textbf{RQ2} & Suy giảm theo \textbf{độ dài} ngữ cảnh / đáp án có khác nhau? & tách nhóm theo độ dài \\
\textbf{RQ3} & Mô hình \textbf{thật sự} nhận ra câu không có đáp án, hay chỉ thiên lệch về từ chối? & mức sàn ``luôn từ chối'' · precision/recall từ chối \\
\textbf{RQ4} & Transformer có vượt baseline TF-IDF ở \textbf{mọi} nhóm? & so từng nhóm \\
\bottomrule
\end{tabularx}

\vspace{0.8em}
\textbf{Giả thuyết ghi trước (RQ1):} PhoBERT yếu hơn ở ranh giới từ ghép vì lỗi tách từ lan truyền vào mô hình.
\end{frame}

% ─────────────────────────────────────────────────────────────── 2 PHƯƠNG PHÁP
\section{Phương pháp}

\begin{frame}{Giao thức: validation chỉ được chấm \emph{một lần}}
\begin{columns}[T]
\begin{column}{0.52\textwidth}
\footnotesize
\begin{tabular}{@{}l r l@{}}
\toprule
\textbf{Tập} & \textbf{Câu} & \textbf{Vai trò} \\
\midrule
fit (train) & \NFit & huấn luyện \\
dev (5\% đoạn văn train) & \NDev & chọn epoch \\
\textbf{validation} & \textbf{\NVal} & \textbf{test, chấm 1 lần} \\
test chính thức & 7.301 & \textcolor{badred}{ẩn nhãn} \\
\bottomrule
\end{tabular}

\vspace{0.6em}
$\bullet$ Chia theo \emph{đoạn văn}; \code{assert} không đoạn văn nào dùng chung\\
$\bullet$ Stress-test \emph{chỉ để đánh giá} (khung mã cũ trộn 60\% vào train)\\
$\bullet$ 1 seed (42), không tìm siêu tham số
\end{column}
\begin{column}{0.44\textwidth}
\small
\textbf{Bốn hệ thống, cùng một harness}
\vspace{0.3em}

\footnotesize
\begin{tabular}{@{}l l@{}}
\textcolor{muted}{\textbf{Luôn từ chối}} & mức sàn \\
\textcolor{tfidf}{\textbf{TF-IDF}} & truy hồi câu (Buổi 4) \\
\textcolor{xlmr}{\textbf{XLM-R-base}} & fine-tune, 384 token \\
\textcolor{phobert}{\textbf{PhoBERT-base-v2}} & fine-tune, 256 token \\
\end{tabular}

\vspace{0.6em}
Cùng vòng lặp, batch 24, 3 epoch, đáp án tối đa 64 token. PhoBERT \emph{buộc} dùng 256 token; 384 token và độ chồng lấp cửa sổ của XLM-R là \emph{lựa chọn} của nhóm, không phải ràng buộc.
\end{column}
\end{columns}
\end{frame}

\begin{frame}{Đưa PhoBERT vào bài toán trích xuất --- không cần Java}
\begin{columns}[T]
\begin{column}{0.55\textwidth}
\small
PhoBERT chỉ có \emph{slow tokenizer} $\Rightarrow$ \textbf{không có offset} (vị trí ký tự).\\
Không có offset thì không cắt được đáp án từ văn bản gốc.

\vspace{0.5em}
\textbf{Lớp \code{SegmentedTokenizer}:}
\begin{enumerate}\footnotesize
\item tách từ bằng \code{pyvi}
\item căn mỗi từ về khoảng ký tự \emph{trong văn bản gốc}
\item BPE từng từ; mọi mảnh mang chung khoảng của từ
\end{enumerate}

\vspace{0.3em}
\footnotesize Kiểm thử: giải mã nhãn ra lại đúng đáp án vàng; 25/25 dự đoán là chuỗi con sạch của ngữ cảnh.
\end{column}
\begin{column}{0.42\textwidth}
\begin{block}{Lỗi ``sai âm thầm'' đã tránh}
\footnotesize
\code{decode()} trả về \code{"Hà\_Nội"}\\
chuẩn hoá xoá dấu câu $\Rightarrow$ \code{"hànội"}\\
so với vàng \code{"hà nội"} $\Rightarrow$ \textbf{EM 0, F1 0}\\[3pt]
Một đáp án \emph{đúng} bị chấm sai --- khung mã cũ mắc lỗi này.
\end{block}
\end{column}
\end{columns}
\end{frame}

% ─────────────────────────────────────────────────────────────── 3 CÔNG CỤ ĐO
\section{Kiểm định công cụ đo}

\begin{frame}{Bộ stress-test 250 câu của nhóm: kiểm định \emph{trước} khi dùng}
\begin{columns}[T]
\begin{column}{0.56\textwidth}
\includegraphics[width=\textwidth]{slide_stress_audit.png}
\end{column}
\begin{column}{0.42\textwidth}
\small
\bignum{\AuditAnsInCtx/\AuditAns}{phobert}\\
câu có đáp án chứa đáp án \emph{trong} ngữ cảnh

\vspace{0.4em}
\footnotesize
$\bullet$ 250 câu dùng chung \AuditCtx{} đoạn văn\\
$\bullet$ E5: 50 câu, \AuditEFiveDistinctQ{} câu hỏi khác nhau\\
$\bullet$ ``Luôn từ chối'' đạt \textbf{\AbstainStressEM} điểm\\[4pt]
\textbf{Nguyên nhân:} \code{\_find\_answer\_start} trả về 0 khi không tìm thấy; ngữ cảnh gán lại theo vòng.

\vspace{0.3em}
\alert{$\Rightarrow$ Không dùng để xếp hạng mô hình.}
\end{column}
\end{columns}
\src{results/stress\_test\_audit.json}
\end{frame}

\begin{frame}{Trần EM do biên từ --- đo được mà không cần mô hình}
\begin{columns}[T]
\begin{column}{0.4\textwidth}
\small
Trên \SegN{} câu có đáp án của validation:

\vspace{0.4em}
\bignum{\SegMisPct\%}{phobert}\\
đáp án cắt ngang một từ (\SegMis{} câu)

\vspace{0.4em}
\bignum{\SegOracle\%}{xlmr}\\
EM tối đa của PhoBERT \emph{hoàn hảo}
\end{column}
\begin{column}{0.56\textwidth}
\footnotesize
\textbf{Đọc từng câu lệch biên:}\\[3pt]
\renewcommand{\arraystretch}{1.25}
\begin{tabular}{@{}l r@{}}
\toprule
\code{pyvi} dính dấu chấm: \code{CNN.} & \MisPeriod \\
\code{pyvi} nối nhầm: \code{Mỹ\_Chú chuột} & \MisMerge \\
Lỗi gán nhãn: ``o thời gian'' & \MisAnnot \\
\textbf{Ranh giới từ ghép thật}: \code{bởi\_vì} & \textbf{\MisCompound} \\
\bottomrule
\end{tabular}

\vspace{0.5em}
\alert{Kênh ``không trả lời được nửa từ'' lấy đi tối đa \SegCeilingLoss{} điểm EM.}
\end{column}
\end{columns}
\src{results/segmentation\_validation.json, results/misaligned\_labels.json}
\end{frame}

% ─────────────────────────────────────────────────────────────── 4 KẾT QUẢ
\section{Kết quả}

\begin{frame}{Kết quả tổng ở ngưỡng mặc định $\tau = 0$ --- validation, \NVal{} câu}
\begin{columns}[T]
\begin{column}{0.58\textwidth}
\includegraphics[width=\textwidth]{slide_main_results.png}
\end{column}
\begin{column}{0.4\textwidth}
\small
\begin{tabular}{@{}l r r@{}}
\toprule
& \textbf{EM} & \textbf{F1} \\
\midrule
Luôn từ chối & \AbstainEM & \AbstainFone \\
TF-IDF & \TfidfEM & \TfidfFone \\
\textcolor{xlmr}{\textbf{XLM-R}} s42 & \XlmrEM & \XlmrFone \\
\textcolor{xlmr}{\textbf{XLM-R}} s13 & \XlmrThirteenEM & \XlmrThirteenFone \\
\textcolor{phobert}{\textbf{PhoBERT}} s42 & \PhobertEM & \PhobertFone \\
\textcolor{phobert}{\textbf{PhoBERT}} s13 & \PhobertThirteenEM & \PhobertThirteenFone \\
\bottomrule
\end{tabular}

\vspace{0.6em}
\scriptsize PhoBERT $-$ XLM-R (EM, cùng seed):\\ s42 \textbf{\PairDiff} · s13 \textbf{+\PairSDiff}\\[2pt]
\alert{Hai seed, hai câu trả lời khác nhau!}
\end{column}
\end{columns}
\src{results/eval\_*\_validation.json, results/diagnosis\_validation.json}
\end{frame}

\begin{frame}{Seed 42, $\tau = 0$: hai hồ sơ kỹ năng \emph{ngược nhau}?}
\begin{columns}[T]
\begin{column}{0.48\textwidth}
\centering
\textbf{Câu có đáp án} ({\NValAns})\\[4pt]
\bignum{+\PairAnsDiff}{phobert}\\
\small EM cho PhoBERT\\
\footnotesize CI [\PairAnsLo; \PairAnsHi] · \textcolor{phobert}{\PhobertAnsEM} vs \textcolor{xlmr}{\XlmrAnsEM}
\end{column}
\begin{column}{0.48\textwidth}
\centering
\textbf{Câu không có đáp án} ({\NValImp})\\[4pt]
\bignum{\PairImpDiff}{xlmr}\\
\small EM (PhoBERT kém hơn)\\
\footnotesize CI [\PairImpLo; \PairImpHi] · \textcolor{phobert}{\PhobertImpEM} vs \textcolor{xlmr}{\XlmrImpEM}
\end{column}
\end{columns}

\vspace{1.2em}
\begin{block}{Nhưng PhoBERT từ chối \PhobertAbsRate\%, XLM-R từ chối \XlmrAbsRate\%}
\small Hai chênh lệch phần lớn là \emph{một}: ai sẵn lòng trả lời hơn. Seed 13 \emph{không} có hai hồ sơ ngược
nhau --- PhoBERT từ chối \PhobertThirteenAbsRate\% --- nên hình ảnh này là của \emph{một lần chạy}.
\end{block}
\src{results/diagnosis\_validation.json, results/thresholds.json}
\end{frame}

\begin{frame}{So đúng cách: chọn ngưỡng $\tau$ cho từng mô hình trên dev}
\begin{columns}[T]
\begin{column}{0.5\textwidth}
\small $\tau$ chọn trên dev (\NDevFull{} câu tách từ train), áp \textbf{một lần} lên validation\\[4pt]
\footnotesize
\renewcommand{\arraystretch}{1.2}
\begin{tabular}{@{}l r r r@{}}
\toprule
& $\tau$ & \textbf{F1} & \textbf{Từ chối} \\
\midrule
\textcolor{xlmr}{\textbf{XLM-R}} s42 & \ThrXlmrTau & \ThrXlmrFone & \ThrXlmrAbsRate\% \\
\textcolor{xlmr}{\textbf{XLM-R}} s13 & \ThrXlmrThirteenTau & \ThrXlmrThirteenFone & \ThrXlmrThirteenAbsRate\% \\
\textcolor{phobert}{\textbf{PhoBERT}} s42 & \ThrPhobertTau & \ThrPhobertFone & \ThrPhobertAbsRate\% \\
\textcolor{phobert}{\textbf{PhoBERT}} s13 & \ThrPhobertRunDTau & \ThrPhobertRunDFone & \ThrPhobertRunDAbsRate\% \\
\bottomrule
\end{tabular}

\end{column}
\begin{column}{0.46\textwidth}
\centering
\bignum{+\PairTDiff{} / +\PairSTDiff}{phobert}\\
\small EM cho PhoBERT, seed 42 / seed 13\\
\footnotesize CI theo bài [\PairTArtLo; \PairTArtHi] / [\PairSTArtLo; \PairSTArtHi]\\[6pt]
Phần chia theo loại câu \emph{không} ổn định giữa hai seed
\end{column}
\end{columns}

\vspace{0.3em}
\begin{block}{``Hai hồ sơ ngược nhau'' là do một lần chạy lệch hiệu chỉnh}
\footnotesize Ở ngưỡng đúng, PhoBERT tốt hơn về tổng ở \emph{cả hai} seed. Nhưng phần hơn đó là hiệu chỉnh:
xem slide sau.
\end{block}
\src{results/thresholds.json, results/diagnosis\_validation.json}
\end{frame}

\begin{frame}{Vậy PhoBERT có \emph{đọc} tốt hơn? \alert{Không đo được khác biệt.}}
\begin{columns}[T]
\begin{column}{0.52\textwidth}
\small Phải khớp \textbf{cả hai} biến gây nhiễu: cấu hình cửa sổ và lần chạy sụp đổ.\\[4pt]
\footnotesize PhoBERT s13 vs XLM-R \textbf{256/96}, cả hai ổn định, mỗi bên ở $\tau$ của mình:
\vspace{0.4em}
\renewcommand{\arraystretch}{1.25}
\begin{tabular}{@{}l r l@{}}
\toprule
\textbf{Nhóm câu} & \textbf{Chênh EM} & \textbf{CI 95\%} \\
\midrule
Toàn bộ & \textbf{+\PairMDiff} & [\PairMArtLo; \PairMArtHi] \\
Có đáp án & \textbf{\PairMAnsDiff} & [\PairMAnsArtLo; \PairMAnsArtHi] \\
Không có đáp án & \textbf{+\PairMImpDiff} & [\PairMImpArtLo; \PairMImpArtHi] \\
\bottomrule
\end{tabular}
\end{column}
\begin{column}{0.44\textwidth}
\small
$\bullet$ CI chỉ rộng $\pm$2 điểm $\Rightarrow$ \textbf{không khác biệt}, chứ không phải thiếu dữ liệu\\[4pt]
$\bullet$ Toàn bộ phần hơn nằm ở câu \textbf{không có đáp án}\\[4pt]
$\bullet$ Lát cắt ``cả hai cùng trả lời'' \textbf{đổi dấu} theo lần chạy: \BothMGap{} với seed 13,
\,$-$\BothMSGap{}\, với lr $10^{-5}$

\vspace{0.5em}
\alert{Điều PhoBERT làm tốt hơn là biết khi nào \emph{không} trả lời.}
\end{column}
\end{columns}
\src{results/diagnosis\_validation.json (paired\_matched\_windows\_tuned)}
\end{frame}

\begin{frame}{RQ1 --- Word-level có yếu ở ranh giới từ? \alert{Không.}}
\begin{columns}[T]
\begin{column}{0.5\textwidth}
\small
\renewcommand{\arraystretch}{1.25}
\begin{tabular}{@{}l r r@{}}
\toprule
& \textcolor{xlmr}{\textbf{XLM-R}} & \textcolor{phobert}{\textbf{PhoBERT}} \\
\midrule
EM, khớp biên (\SegAligned) & \XlmrAlEM & \textbf{\PhobertAlEM} \\
EM, lệch biên (\SegMis) & \textbf{\XlmrMisEM} & \PhobertMisEM \\
F1, lệch biên & \XlmrMisFone & \textbf{\PhobertMisFone} \\
\midrule
Lỗi biên / lần trả lời & \XlmrBoundaryRate\% & \PhobertBoundaryRate\% \\
\bottomrule
\end{tabular}
\end{column}
\begin{column}{0.46\textwidth}
\small
$\bullet$ Cơ chế \textbf{có thật}: câu lệch biên, PhoBERT F1 cao mà EM thấp\\
\footnotesize\quad vàng ``trận Gonzales'' $\to$ \code{tại\_trận} $\to$ ``tại trận Gonzales''\\[4pt]
\small $\bullet$ Nhưng chỉ \SegMis{} câu, $p$ = \PairMisP, chặn ở \SegCeilingLoss{} điểm EM\\[4pt]
$\bullet$ Khi đã trả lời, tỉ lệ lỗi biên \textbf{như nhau}

\vspace{0.5em}
\alert{Giả thuyết ghi trước sai hướng.}
\end{column}
\end{columns}
\src{results/diagnosis\_validation.json, results/segmentation\_validation.json}
\end{frame}

\begin{frame}{RQ3 --- Nhận ra câu không có đáp án? \alert{Một phần.}}
\begin{columns}[T]
\begin{column}{0.5\textwidth}
\small
\renewcommand{\arraystretch}{1.25}
\begin{tabular}{@{}l r r r@{}}
\toprule
& \textbf{Từ chối} & \textbf{Prec.} & \textbf{Recall} \\
\midrule
Luôn từ chối & \AbstainAbsRate & \AbstainAbsPrec & \AbstainAbsRec \\
\textcolor{xlmr}{\textbf{XLM-R}} & \XlmrAbsRate & \XlmrAbsPrec & \XlmrAbsRec \\
\textcolor{phobert}{\textbf{PhoBERT}} & \PhobertAbsRate & \PhobertAbsPrec & \PhobertAbsRec \\
\bottomrule
\end{tabular}

\vspace{0.4em}\footnotesize Tỉ lệ câu không có đáp án thật: \PctValImp\%
\end{column}
\begin{column}{0.46\textwidth}
\small
$\bullet$ \textcolor{xlmr}{\textbf{XLM-R}}: từ chối đúng \emph{tần suất}, sai \emph{đối tượng} (\XlmrErrFalseAbs{} câu có đáp án bị từ chối)\\[4pt]
$\bullet$ \textcolor{phobert}{\textbf{PhoBERT}}: gần như không từ chối\\[4pt]
$\bullet$ Rơi vào \textbf{bẫy} của người gán nhãn: câu trả lời sai trùng đúng cụm ``plausible'' ---
PhoBERT \PhobertTrapPct\%, XLM-R \XlmrTrapPct\%
\end{column}
\end{columns}
\src{results/diagnosis\_validation.json}
\end{frame}

\begin{frame}{Vì sao PhoBERT không từ chối? Một sự kiện khi huấn luyện}
\begin{columns}[T]
\begin{column}{0.5\textwidth}
\small Huấn luyện tiếp từ epoch 1, \textbf{cùng batch}, chỉ khác lr:\\[3pt]
\scriptsize
\renewcommand{\arraystretch}{1.1}
\begin{tabular}{@{}l r r r r@{}}
\toprule
& \multicolumn{2}{c}{\textbf{lr $2{,}2\cdot10^{-5}$}} & \multicolumn{2}{c}{\textbf{lr $10^{-5}$}} \\
\textbf{Bước} & \textbf{loss} & \textbf{grad} & \textbf{loss} & \textbf{grad} \\
\midrule
\tabinput{../report/generated/tab_probe_resume_slide.tex}
\bottomrule
\end{tabular}\\[2pt]
loss: ``không có đáp án'' (TB khối 100 bước); grad: chuẩn lớn nhất
\end{column}
\begin{column}{0.46\textwidth}
\small
$\bullet$ Dữ liệu hỏng? \textbf{Không}: \ScanPhobertViolations{} vi phạm / \ScanPhobertFeatures{} feature\\[3pt]
$\bullet$ Tính chất của PhoBERT? \textbf{Không}: epoch 1 từ chối \DevEpochOneAbsRate\% (dev)\\[3pt]
$\bullet$ Bất ổn tối ưu hoá: gradient vọt \textbf{\ResumeHighMaxSpike}, loss ``không đ.án'' kẹt ở \ResumeHighClsEnd; lr $10^{-5}$ thì không\\[3pt]
$\bullet$ Khả năng phân biệt \textbf{còn}: điểm \code{[CLS]} chỉ dịch $\approx$ \ThrPhobertTau{} logit\\[3pt]
$\bullet$ \textbf{Phụ thuộc lần chạy}: seed 13, cùng lr $3\cdot10^{-5}$, không sụp đổ
\end{column}
\end{columns}
\src{results/probe\_resume.json, feature\_scan\_phobert.json}
\end{frame}

\begin{frame}{Mọi lần chạy trên dev: $\tau = 0$ và $\tau$ hiệu chỉnh}
\small F1 trên toàn bộ dev (\NDevFull{} câu tách từ train) --- dữ liệu dùng cho mọi lựa chọn.\\[4pt]
\footnotesize
\renewcommand{\arraystretch}{1.2}
\begin{tabular}{@{}l l r r r r r@{}}
\toprule
& & \multicolumn{2}{c}{\textbf{$\tau = 0$}} & \multicolumn{3}{c}{\textbf{$\tau$ chọn trên dev}} \\
\textbf{Lần chạy} & & \textbf{F1} & \textbf{Từ chối} & $\tau$ & \textbf{F1} & \textbf{Từ chối} \\
\midrule
\tabinput{../report/generated/tab_dev_runs.tex}
\bottomrule
\end{tabular}

\vspace{0.6em}
\small \alert{Hai lần chạy sụp đổ đều là seed 42.} Sau hiệu chỉnh, hai seed PhoBERT gần như trùng nhau; lr $10^{-5}$ ổn định nhưng còn thiếu epoch. Nhóm báo cáo \textbf{hai seed ngang nhau}.
\src{results/thresholds.json, results/training\_curve\_*.json}
\end{frame}

\begin{frame}{RQ2 độ dài · RQ4 baseline · và bộ stress-test}
\begin{columns}[T]
\begin{column}{0.5\textwidth}
\includegraphics[width=\textwidth]{slide_by_answer_length.png}
\footnotesize \textbf{RQ2:} không suy giảm theo độ dài ngữ cảnh; đáp án dài làm giảm EM (F1 giữ nguyên) như nhau ở cả hai.\\[3pt]
\textbf{RQ4:} Transformer vượt TF-IDF ở mọi nhóm; TF-IDF (EM \TfidfEM) còn thua ``luôn từ chối''.
\end{column}
\begin{column}{0.46\textwidth}
\small \textbf{EM trên stress-test 250 câu}\\[3pt]
\begin{tabular}{@{}l r@{}}
\toprule
\textbf{Luôn từ chối} & \textbf{\AbstainStressEM} \\
XLM-R & \XlmrStressEM \\
PhoBERT & \PhobertStressEM \\
TF-IDF & \TfidfStressEM \\
\bottomrule
\end{tabular}

\vspace{0.5em}
\alert{``Không làm gì'' đứng đầu} --- đúng như kiểm định dự báo. Bộ này không dùng để xếp hạng.
\end{column}
\end{columns}
\end{frame}

% ─────────────────────────────────────────────────────────────── 5 KẾT LUẬN
\section{Kết luận}

\begin{frame}{Kết luận}
\small
\begin{enumerate}
\item \textbf{So ở $\tau$ cố định với một seed là so hiệu chỉnh, không phải năng lực.} Seed 42, $\tau = 0$: ``bằng nhau, hồ sơ ngược nhau''. Ở $\tau$ chọn trên dev: PhoBERT hơn \PairTDiff{} / \PairSTDiff{} điểm EM ở hai seed.
\item \textbf{PhoBERT seed 42 ``không từ chối'' là một sự kiện khi huấn luyện}, tái hiện được, kèm gradient vọt; seed 13 không gặp.
\item \textbf{Không đo được khác biệt về ``đọc''}: khớp cửa sổ thì chênh lệch trên câu có đáp án là \PairMAnsDiff{} điểm. Tách từ cũng không quyết định: chỉ \SegMisPct\% đáp án cắt ngang từ.
\item \textbf{Kiểm định công cụ đo trước khi dùng}: bộ stress-test của chính nhóm xếp ``không làm gì'' đứng đầu.
\end{enumerate}

\vspace{0.4em}
\begin{block}{Giới hạn}
\footnotesize 2 seed · 3 epoch, không tìm siêu tham số · \code{pyvi} thay RDRSegmenter · chưa rõ vì sao sụp đổ ở seed 42 mà không ở seed 13 · validation = public test VLSP (private test ẩn nhãn)
\end{block}
\end{frame}

\begin{frame}{Hướng phát triển}
\small
\begin{itemize}
\item Vì sao sự kiện sụp đổ chỉ xảy ra với PhoBERT; có lặp lại trên CUDA không; warmup dài hơn thay vì hạ lr
\item Huấn luyện thêm epoch (XLM-R và PhoBERT lr $10^{-5}$ còn đang cải thiện); từ 3 seed trở lên
\item PhoBERT với RDRSegmenter; sửa luật dính dấu chấm của \code{pyvi}
\item Xây lại bộ stress-test có kiểm định tự động và câu hỏi do người viết
\end{itemize}

\vspace{1em}
\centering
{\headfont\bfseries\Large Cảm ơn thầy và các bạn đã lắng nghe!}\\[4pt]
\footnotesize Mã nguồn, kết quả và nhật ký: \code{github.com/minhdang900/uit-benchmark-transformer-model-mrc}
\end{frame}

\appendix
\begin{frame}{Phụ lục --- Quá trình huấn luyện}
\begin{columns}[T]
\begin{column}{0.58\textwidth}
\includegraphics[width=\textwidth]{slide_training_curves.png}
\end{column}
\begin{column}{0.4\textwidth}
\footnotesize
Chọn epoch theo \textbf{F1 trên dev} (tách từ train).\\[3pt]
XLM-R: epoch \XlmrBestEpoch, \XlmrTrainMin{} phút.\\
PhoBERT lr $3\cdot10^{-5}$: epoch \PhobertRunABestEpoch, \PhobertRunATrainMin{} phút.\\
PhoBERT lr $2\cdot10^{-5}$: epoch \PhobertRunBBestEpoch.\\
PhoBERT lr $10^{-5}$: epoch \PhobertRunCBestEpoch.\\
PhoBERT seed 13: epoch \PhobertRunDBestEpoch.\\[3pt]
Hai seed báo cáo ngang nhau.
\end{column}
\end{columns}
\end{frame}

\begin{frame}{Phụ lục --- Độ dao động theo seed và cửa sổ}
\small Validation, EM / F1. Seed tách dev cố định (42) để điểm dev so sánh được.\\[4pt]
\footnotesize
\renewcommand{\arraystretch}{1.2}
\begin{tabular}{@{}l r l r r@{}}
\toprule
\textbf{Mô hình} & \textbf{Seed} & \textbf{Cửa sổ} & \textbf{$\tau = 0$} & \textbf{$\tau$ dev} \\
\midrule
\tabinput{../report/generated/tab_seeds.tex}
\bottomrule
\end{tabular}
\src{results/eval\_*\_validation.json, results/thresholds.json}
\end{frame}

